\subsection{A rational quadratic-Hilbert consequence of the Euler-system machinery}
\label{subsec:lm-lz-rational}

In this subsection LZ denotes the February 2025 version of
\citet{LoefflerZerbes2025}, with theorem numbers and pages referring
to that version.  The distinction between its
stated large-image theorem and its formal intermediate results is
essential.  Its condition (BI), Definition 2.5.1, requires a product of
two special linear groups and fails for CM representations.  The next
proposition is instead a deduction from the constructions in Sections
3.4--3.5 and 6--8 and the abstract results in Sections 9--10.  Only a
rational Selmer conclusion is sought; the canonical integral period
normalization of Section 3.2.3 is not required.

This is an arithmetic input for the proper finite-descent application
at exponent thirteen. The regular-prime quotient-section classification
proved earlier does not depend on this separate source chain.
The following table fixes the versions and roles of its inputs.

\begin{center}
\begingroup
\small
\setlength{\tabcolsep}{3pt}
\renewcommand{\arraystretch}{1.12}
\begin{tabular}{@{}p{.17\linewidth}p{.17\linewidth}p{.30\linewidth}p{.30\linewidth}@{}}
\toprule
Source & Cited version & Results used & Role in the argument \\
\midrule
\citet{MazurRubin2004}
& Published memoir, 2004
& \S3.5, pp.~27--28: hypotheses (H.0)--(H.4)
& The lattice and auxiliary-twist hypotheses checked directly below. \\
\addlinespace
\citet{LoefflerRockwoodZerbes2024}
& Author version v3, 30 July 2024
& Theorems~2.6.2 and~7.1.1, pp.~4 and~15
& Finite generation and ordinary families of cohomology classes, used in the family Euler-system construction. \\
\addlinespace
\citet{LoefflerZerbes2025}
& Author version v2, 18 February 2025
& Theorems~3.4.4, 3.5.1 and~7.1.3; \S7.2; Theorem~8.1.1, p.~22
& The Hilbert family, its Yoshida lift and $p$-adic $L$-functions, and uniformly integral family Euler classes. \\
\addlinespace
\citet{LoefflerZerbes2026}
& Author version v4, 21 April 2026
& Theorem~A, p.~4; Theorem~6.6.4, p.~14
& Explicit reciprocity in the range $r_1-r_2\ge3$, followed here by continuation from a dense set of such weights. \\
\addlinespace
\citet{LoefflerZerbes2025}
& The same February 2025 version
& Theorem~8.2.1; Lemma~10.1.1; Propositions~10.1.4 and~10.2.2; Theorem~10.3.3; Corollary~9.6.2; proof of Theorem~12.3.1, p.~33
& Meromorphic comparison, leading terms and characteristic-zero Selmer control. These intermediate results supply the rational deduction; the stated (BI) theorem is not applied. \\
\bottomrule
\end{tabular}
\endgroup
\end{center}

\begin{proposition}[Rational quadratic-Hilbert criterion]
\label{prop:lm-rational-hilbert}
Let $F/\mathbb Q$ be real quadratic, and let $\pi$ be a cuspidal Hilbert
automorphic representation of weight $(2,2,0,0)$.  Suppose that its
finite central character factors through $N_{F/\mathbb Q}$ and that
$\pi^\sigma\not\simeq\pi$, where $\sigma$ is the nonidentity automorphism
of $F$.  Let $p\geq7$ split in $F$, be prime to the level, and be ordinary
at both places above $p$.  Write
\[
 V=\operatorname{Ind}_{G_F}^{G_{\mathbb Q}}\rho_{\pi,p}^{\vee}.
\]
Assume the following conditions for a stable lattice $T\subset V$.
The hypotheses (H.0)--(H.4) of
\citet[\S3.5, pp.~27--28]{MazurRubin2004} hold for
$T$ and for $T(\eta)$ whenever $\eta$ has $p$-power order and square-free
conductor prime to $p$ and the level and discriminant in this setup.
Assume also that $V$ is irreducible.
Then
\[
 L(\pi,1)\ne0\quad\Longrightarrow\quad H^1_f(F,\rho_{\pi,p}^{\vee})=0.
\]
\end{proposition}

\begin{proof}
We explain the deduction, including the hypotheses that are sometimes
hidden by citing the more restrictive main theorem.  Choose $S$ to
contain $p$ and all bad primes, and let $\mathcal R$ consist of
square-free products of primes $q\notin S$ with $q\equiv1\pmod p$.
The group $\Delta_m$ is the maximal $p$-power quotient of
$(\mathbb Z/m\mathbb Z)^\times$.  Thus the auxiliary characters in LZ
Section 9.1 are exactly of the type stipulated here.  They are not
arbitrary finite-order Dirichlet characters.

By LZ Theorem 3.4.4, there is an ordinary Hilbert family on a weight
polydisc through $\pi$.  Theorems 3.5.1 and 7.1.3 give its analytic
$p$-adic $L$-functions and twisted Yoshida lift.  Restrict to the
one-parameter family of Section 7.2, of weights
\[
 (2+4\mu,2+2\mu,-2\mu,-\mu).
\]
Let $\mathcal M_{\mathrm{et}}$ be the torsion-free pullback of the
Siegel Galois sheaf, a free rank-four module on this disc, and use
\[
 \mathcal M=\mathcal M_{\mathrm{et}}^{\vee}(\mathbf t_1)
\]
as in LZ Section 11.  Here $\mathbf t_1$ is the family cyclotomic
twist whose value at a classical point is $t_1$.  Proposition 7.2.4
identifies the semisimplification of the fiber of
$\mathcal M_{\mathrm{et}}$ with
$\operatorname{Ind}\rho_{\pi,p}(t_1)$; hence the fiber of
$\mathcal M$ at the central point, where $t_1=0$, has
semisimplification $V$.  The assumed irreducibility identifies the
representation itself.
The Hilbert family construction requires no smoothness of the entire
Siegel eigenvariety at this point.

LZ Theorem~8.1.1, p.~22, supplies a family of Euler-system classes
for this module. Its construction uses
\citet[Theorems~2.6.2 and~7.1.1, pp.~4 and~15]{LoefflerRockwoodZerbes2024},
which give finite generation and the ordinary family classes without
a large-image hypothesis.
Its last assertion places all the classes in a single finitely generated
module over the power-bounded functions on the disc.  This uniform
integrality assertion is prior to, and independent of, condition (BI).
There is no lattice mismatch in applying the hypotheses.  If one stable
lattice in an irreducible representation has absolutely irreducible
reduction, every stable lattice is homothetic to it: scale the second
lattice into the first but not into its maximal-ideal multiple.  Its
nonzero residual image is the whole irreducible reduction, and
Nakayama's lemma makes the inclusion an equality.  Thus the lattice
obtained from the family has the same Mazur--Rubin properties.  The
same argument applies after any of the allowed coefficient extensions
and auxiliary twists.
Fix the smoothing parameter $r=3$, and an integer $c>1$ prime to the
bad primes, with the finite central character taking value $1$ at $c$.
Enlarge $S$ by the prime factors of $c$ and use the corresponding
$\mathcal R$, so that every Euler-system index is prime to $c$.
For the sign containing the central critical value, Theorem 8.2.1
expresses the regulator of these classes as the product of the desired
family $p$-adic $L$-function, a smoothing factor, and an auxiliary
one-variable family $L$-value.  Here the completed explicit reciprocity
law is used in its current scope:
\citet[Theorem~A, p.~4; Theorem~6.6.4, p.~14 of the April 2026
author version]{LoefflerZerbes2026} assumes $r_1-r_2\geq3$.  For our family this difference
is $k_2-2=2\mu$.  The classical weights $\mu\geq2$ in the weight
disc form a Zariski-dense subset, satisfy this restriction, and also
have $r=3\leq2\mu$.  Thus the analytic-continuation argument proving
LZ Theorem~8.2.1 applies on that dense subset, including its
specialization at $\mu=0$; no small-weight reciprocity statement is
being assumed.  The auxiliary factor is not identically
zero: at sufficiently large classical $\mu$, its evaluation is at the
noncentral critical integer $-\mu+3$, and nonvanishing follows from
Proposition 2.3.7.  It may vanish at $\mu=0$, which is precisely why the
leading-term argument is needed.

The independence of the comparison ratio from the cyclotomic variable
can be checked explicitly in this continuation argument. On a fixed
cyclotomic component choose a nonzero integer $b$ divisible by $p-1$
and prime to $p$. At classical weight $\mu$, the allowed integral
cyclotomic exponents form $[-2\mu,-\mu-1]$. The pairs $(\mu,j)$ for
which both $j$ and $j+b$ lie in that interval are dense in the product
of the two discs: choose arbitrarily large positive $\mu$ in any
prescribed residue class, then choose $j$ in its growing interval in
any prescribed residue class. The classical comparison ratio depends
only on $\mu$, so its meromorphic continuations at $j+b$ and at $j$
agree on this dense set and hence identically. Integer multiples of
$b$ are dense in the additive $p$-adic parameter disc. Continuity,
or differentiation along translations tending to zero, therefore
makes the ratio independent of the cyclotomic variable. This is the
one-variable meromorphic comparison used in LZ Theorem~8.2.1; its
zeros or poles at the central weight are handled by the following
leading-term argument.

Let $X$ be a parameter at the central point.  Divide the entire family
of Euler classes by its largest common power $X^h$.  Lemma 10.1.1 and
Proposition 10.1.4 show that the specialized classes form a nonzero
Euler system in one fixed stable lattice.  In particular no separate
denominator depending on $m\in\mathcal R$ is introduced.  The two nested
ordinary local submodules have ranks two and three, as in Section 11.
Complex conjugation has two eigenvalues of each sign.  The local
vanishing conditions of Assumption 9.5.1 hold by Proposition 3.3.1,
also after the finite unramified extensions at $p$ supplied by $\Delta_m$.
One can equivalently check this by purity: the unramified part of each
ordinary one-dimensional character has a Frobenius eigenvalue of
nonzero Weil weight, so no power of it is $1$.  The same applies to
the dual local subquotients.

The precise hypotheses of Proposition 10.2.2 are now satisfied: they
are (H.0)--(H.4) for $T$ and its $\Delta_m$-twists, not (BI).  Thus the
leading-term Euler system has a nonzero regulator in its rank-one
local quotient.  Divide the meromorphic comparison isomorphism by
the auxiliary family $L$-value and apply Theorem 10.3.3.  It identifies
this regulator with the desired $p$-adic $L$-function times the
smoothing factor, up to a nonzero constant resulting from a basis of
the one-dimensional crystalline space.

At the central critical character the smoothing factor for $r=3$ is
\[
 (c^2-c^4)(c^2-c^{-2})\ne0.
\]
The ordinary interpolation factors are nonzero by Proposition 3.1.5.
Consequently $L(\pi,1)\ne0$ makes the regulator nonzero at that
character.  Corollary 9.6.2 with $k=0$, localized at this characteristic
zero height-one prime of the Iwasawa algebra, says that the ordinary
Selmer characteristic ideal is a unit there.  Multiplying the chosen
lattice or regulator by a power of $p$ does not affect this conclusion.
This is why the integral Betti lattice and the (BI)-dependent
normalization of Section 3.2.3 are unnecessary.

For completeness, let $\Lambda$ denote the cyclotomic Iwasawa algebra
on the component containing the central character, and let $\mathfrak q$
be its characteristic-zero height-one prime. Write
$\widetilde H^i_{\mathrm{Iw}}$ for the cohomology of the Selmer complex
with ordinary conditions at both places above $p$. The preceding
characteristic-ideal conclusion gives
$(\widetilde H^2_{\mathrm{Iw}})_{\mathfrak q}=0$.
By LZ Proposition~3.3.5, p.~11, its first cohomology injects into global
Iwasawa cohomology and is torsion-free; the Euler-characteristic rank
formula gives rank zero, so $\widetilde H^1_{\mathrm{Iw}}=0$ on this
component. The derived base-change argument in the proof of LZ
Theorem~12.3.1, p.~33, now gives
\[
0\longrightarrow
\widetilde H^1_{\mathrm{Iw}}\otimes_\Lambda\kappa(\mathfrak q)
\longrightarrow \widetilde H^1(\text{central fiber})
\longrightarrow
\operatorname{Tor}_1^\Lambda
 (\widetilde H^2_{\mathrm{Iw}},\kappa(\mathfrak q))
\longrightarrow0.
\]
Both end terms are zero. Global degree-zero and degree-three terms,
and the local exceptional terms, vanish rationally by irreducibility
and the local and dual invariant vanishing established above.
At this weight-two critical character the ordinary subspace is a
cyclotomic twist of an unramified line and its quotient is unramified; their nontrivial
Frobenius eigenvalues identify the ordinary local condition with the
Bloch--Kato condition. Thus the central-fiber group is the required
finite Selmer group, which vanishes. This is the formal base-change
deduction, rather than an application of the (BI) version of the stated
main theorem. The integral period normalization is unnecessary because
powers of $p$ are units at $\mathfrak q$.
\end{proof}

\subsection{The cubic twists of the cyclic-quartic CM surface}

Let $\mathcal K$ be the Klein quintic
\[
 A^4B+B^4C+C^4A=0,
\]
let $c$ be the rational order-three coordinate permutation, and let
$B=\mathcal K/\langle c\rangle$.  Thus $B$ has the genus-two model
\[
 v^2=1-12t+52t^2-104t^3+104t^4-52t^5.
\]
Put $A_B=\operatorname{Jac}(B)$, $L=\mathbb Q(\zeta_{13})$,
$H=\langle3\rangle\subset(\mathbb Z/13\mathbb Z)^\times$, and
$K=L^H$.  Then $K$ is cyclic quartic CM, with real quadratic subfield
$F=\mathbb Q(\sqrt{13})$.  Let $\chi$ be a nontrivial cubic character
modulo $13$.

\begin{lemma}[The CM algebra is defined over the quartic field]
\label{lem:lm-thirteen-cm-field}
There is a degree-four CM algebra
$K\hookrightarrow\operatorname{End}^0_K(A_B)$.  On the four CM lines
the action of $G_{\mathbb Q}$ is the natural transitive cyclic
permutation through $\operatorname{Gal}(K/\mathbb Q)$.
\end{lemma}
\begin{proof}
Let $\tau$ denote the order-thirteen diagonal automorphism of
$\mathcal K$ over $L$, with $c\tau c^{-1}=\tau^3$.  The operator
\[
 \theta=\tau+\tau^3+\tau^9
\]
commutes with $c$ and hence induces an element of
$\operatorname{End}^0_{\overline{\mathbb Q}}(A_B)$; explicitly use
pullback, this operator, and norm,
and divide by the degree $3$ of the quotient map.  The Galois
subgroup $H$ fixes this operator, since it permutes its three terms.
The induced endomorphism is consequently defined over $K$.

The quotient by $\tau$ is the projective line and its three branch
points are totally ramified.  Equivalently each nonidentity power of
$\tau$ has three fixed points, so its trace on first cohomology is
$-1$.  It follows that the nontrivial order-thirteen characters
occur once each in the Tate module of
$\operatorname{Jac}(\mathcal K)$.  Taking the
$c$-invariant part leaves one line for each of the four $H$-orbits.
The eigenvalues of $\theta$ are the Gaussian periods
\[
 \sum_{h\in H}\zeta_{13}^{ah},
 \qquad aH\in(\mathbb Z/13\mathbb Z)^\times/H.
\]
They are distinct: equality of two such sums would give a nonzero
rational polynomial of degree at most $12$, with constant coefficient
zero, divisible by $1+X+\cdots+X^{12}$.  They are therefore the four
conjugates of a generator of $K$.  This proves the asserted embedding.
Galois conjugation takes the period indexed by $aH$ to the period
indexed by $baH$, proving the permutation assertion.  Only this CM
subalgebra is needed below; no larger endomorphism-field assertion
is being assumed.
\end{proof}

\begin{lemma}[Explicit Mazur--Rubin conditions]
\label{lem:lm-thirteen-mr}
For all but finitely many primes $p\geq7$ splitting completely in
$K\mathbb Q(\zeta_3)$, each of the representations
\[
 V_p(A_B)\otimes\chi^a,\qquad a=1,2,
\]
and their duals and Tate twists satisfies (H.0)--(H.4).  The same is
true after a further twist by any $p$-power-order Dirichlet character
$\eta$ whose conductor is prime to $13p$.  In particular it is true
for the square-free auxiliary characters of
Proposition~\ref{prop:lm-rational-hilbert}.
\end{lemma}
\begin{proof}
Discard the finitely many primes dividing the order of the CM
endomorphism lattice, the polarization degree, or the bad reduction
conductor, and choose the four CM idempotents over $\mathbb Z_p$.
The Tate lattice is their direct sum, with each summand rank one.
Index them by $0,1,2,3$ cyclically, with complex conjugation adding
$2$.  Since the cyclic quartic field has no proper CM subfield, its
CM type is, after cyclic relabeling, $\{0,1\}$.

At a split good prime $p$, each CM line is a crystalline character
of $G_{\mathbb Q_p}$ with cyclotomic exponent $0$ or $1$. Restricted
to inertia it is therefore respectively $1$ or the cyclotomic
character $\varepsilon_p$; the remaining crystalline rank-one
factor is unramified.  The cyclotomic character maps inertia onto
$\mathbb Z_p^\times$.  The four Galois conjugates of this inertia
group thus supply, for arbitrary units $u_i$, the four diagonal
matrices
\[
 \begin{split}
 &\operatorname{diag}(u_1,u_1,1,1),\quad
 \operatorname{diag}(1,u_2,u_2,1),\\
 &\operatorname{diag}(1,1,u_3,u_3),\quad
 \operatorname{diag}(u_4,1,1,u_4).
 \end{split}
\]
They lie in $G_K$; both $\chi$ and every allowed $\eta$ are trivial
on them.  Their product is
\[
 \operatorname{diag}(u_1u_4,u_1u_2,u_2u_3,u_3u_4),
 \qquad \varepsilon_p= u_1u_2u_3u_4.
\]
In particular every integral torus element
\[
 \operatorname{diag}(a,b,c,d),\qquad ac=bd,
\]
is supplied: take $(u_1,u_2,u_3,u_4)=(1,b,c/b,a)$.
Its cyclotomic multiplier is $ac$.  This proves the required exact
torus-image assertion by local crystalline characters, without an
unverified index assertion for a global CM reciprocity image.

For (H.1), the four diagonal characters of this torus are pairwise
distinct modulo $p$, and its residual image has order prime to $p$.
Every invariant subspace over an algebraic closure of the residue
field must therefore be a sum of CM lines.  The cyclic Galois action
on these lines is transitive, so there is no proper such subspace.
This is absolute irreducibility.  Scalar twisting by $\chi^a$ does
not change it, and the residual character of $\eta$ is trivial.

For (H.2), choose $g\in I_{13}\cap G_K$ with $\chi(g)$ a primitive
cube root $\zeta$ in $\mathbb Z_p^\times$.  This is possible because
$13$ is totally ramified in $L$ and the subgroup $H$ fixing $K$
maps nontrivially under $\chi$.  The element $g$ has
$\varepsilon_p(g)=1$, and its action on the CM lines is diagonal.
Polarization implies that its two opposite products are both $1$.
Using the torus just constructed, multiply $g$ by a product of
conjugate $p$-inertia elements so that its untwisted Tate action becomes
\[
 D=\operatorname{diag}(\zeta^{-a},1,\zeta^{a},1).
\]
The adjusting product has cyclotomic multiplier $1$.  The resulting
element $\tau$ fixes $\mu_{p^\infty}$ and has $\eta(\tau)=1$;
on the $\chi^a$-twist its eigenvalues are
\[
 1,\ \zeta^a,\ \zeta^{2a},\ \zeta^a.
\]
Exactly one is $1$, and the other three differ from $1$ by units.
Thus $T/(\tau-1)T$ is free of rank one.  Duals have the same property,
and Tate twists do not alter the action of this $\tau$.

For (H.3), use the same inertia construction for $-I_4$, whose
cyclotomic multiplier is $1$.  Denote the resulting Galois element
by $z$.  It is trivial under $\chi$ and $\eta$, and its image is
central in the joint image of the representation and
$\varepsilon_p$.  It acts by $-1$ on the residual module and its
Cartier dual.  Sah's elementary cocycle identity consequently
annihilates the two required first cohomology groups by $2$.
Explicitly, for a continuous cocycle $f$ and a central $z$,
\[
 (z-1)f(g)=(g-1)f(z).
\]
Here $z-1=-2$ is invertible, so $f$ is the coboundary of
$-f(z)/2$.  Thus both groups vanish.  This also applies to the infinite
joint image by the identical continuous-cocycle calculation.
Hypothesis (H.0) is finite freeness and finite ramification; (H.4)
holds since $p\geq5$.  The argument proves every assertion for the
further $\eta$-twists as well.
\end{proof}

\begin{proposition}[The exponent-thirteen split-prime Selmer consequence]
\label{prop:lm-thirteen-split-prime}
Let $P_{13}$ be the rational fourfold complementary to $A_B$ in the
Klein-quintic Jacobian.  The decomposition in Proposition~\ref{prop:thirteenpartial} is
\[
 V_p(P_{13})\simeq
 (V_p(A_B)\otimes\chi)\oplus(V_p(A_B)\otimes\chi^2)
\]
after adjoining the cubic character values, and that proposition proves $L(P_{13},1)\ne0$.  For all but finitely many rational primes $p$ splitting completely in
$K\mathbb Q(\zeta_3)$,
\[
 H^1_f(\mathbb Q,V_p(P_{13}))=0,
 \qquad T_p\Sha(P_{13}/\mathbb Q)=0.
\]
\end{proposition}
\begin{proof}
The CM algebra over $K$ proved above gives algebraic Hecke characters
with the prescribed CM infinity types and the Tate-module realizations
used here: \citet[\S7, Theorem~10, pp.~511--512; Theorem~11(i),
p.~512]{SerreTate1968} constructs the algebraic idelic character and
identifies its $p$-adic realizations with the Tate characters; the
construction on pp.~513--514 gives the associated complex
Gr\"ossencharacters.
The primitive CM character over $K$ then induces over its real quadratic
subfield $F$ a weight-two cuspidal Hilbert representation $\pi_0$.
This is the quadratic induction theorem of
\citet[Proposition~12.1, p.~206 in the retypeset author version]{JacquetLanglands1970}.
Its norm-factorization obstruction to cuspidality is absent because
the two exponents of our CM character at every complex place differ
by $1$.  Its induced real Weil parameter is the weight-two discrete
series parameter; the relation $|a_w-b_w|=k_v-1$ is also recorded
in \citet[Proposition~B.2.7, pp.~654--655]{Nekovar2012}.
Fix the cohomological convention of LZ: $\rho_{\pi_0,p}$ is the
cohomological representation, so its dual induces to
$W=V_p(A_B)$.  Write $\chi$ also for the automorphic character
corresponding to the chosen Galois character in this convention;
explicitly its unramified local value is the value of the Galois
character on the geometric Frobenius used in LZ's Euler polynomial.
Thus automorphic twisting by this $\chi$ twists $\rho_{\pi_0,p}$
by the Galois character $\chi$.  Define
\[
 \pi_a=\pi_0\otimes\chi^{-a},\qquad
 \operatorname{Ind}_{G_F}^{G_{\mathbb Q}}\rho_{\pi_a,p}^{\vee}
      =W\otimes\chi^a,\qquad a=1,2.
\]
This fixes the normalization rather than choosing it after applying
the theorem.  Labeling Dirichlet characters by their values on
arithmetic instead of geometric Frobenius interchanges $a=1,2$
and does not change the pair.  Pairing opposite CM lines by the
polarization shows that its determinant on the Tate side is
cyclotomic.  The finite central character of $\pi_a$ is
$\chi^{-2a}\circ N_{F/\mathbb Q}$, which satisfies the
norm-factorization hypothesis of
Proposition~\ref{prop:lm-rational-hilbert}; it need not be trivial.
The two conjugate Hilbert constituents are distinct because the
four CM characters have distinct infinity types.  Their induction
to $G_{\mathbb Q}$ is irreducible by the same CM-line argument as
in Lemma~\ref{lem:lm-thirteen-mr}.

At the chosen split primes both Hilbert ordinary conditions hold:
the rank-one CM characters are crystalline, and each opposite pair
has one exponent $0$ and one exponent $1$.  All local invariants
required in the Euler-system argument vanish, also after the
unramified $p$-power auxiliary extensions at $p$: a Frobenius power
cannot turn a Weil eigenvalue of absolute value $p^{1/2}$ or
$p^{-1/2}$ into $1$.  Lemma~\ref{lem:lm-thirteen-mr} gives exactly
the auxiliary-twist conditions.  The product formula for the
$L$-function and $L(P_{13},1)\ne0$ imply the nonvanishing of each
of the two twisted central values.  Apply
Proposition~\ref{prop:lm-rational-hilbert} to the two Hilbert forms,
$\pi_1$ and $\pi_2$, with the displayed dual identification.
Shapiro's lemma then
gives finite-Selmer vanishing over $\mathbb Q$ for both summands.
Finite extension of the coefficient field is faithfully flat,
so the asserted vanishing for $V_p(P_{13})$ follows.

The Kummer exact sequence identifies the rational Bloch--Kato
Selmer group of an abelian variety with its pro-$p$ Selmer group
tensored with $\mathbb Q_p$.  Thus the Tate module of its
Tate--Shafarevich group becomes zero after tensoring with
$\mathbb Q_p$.  That Tate module is torsion-free over $\mathbb Z_p$,
so it is zero.  This proves the claim.  It does not claim finiteness
of the entire Tate--Shafarevich group.
\end{proof}

